% KDD 2026 Research Track Submission
% Title: What Makes LLM Structured Outputs Inconsistent? A Comprehensive Factor Analysis
%
% This paper focuses on FACTOR ANALYSIS (different from ICML which focuses on STED method)

\documentclass[sigconf,anonymous,review]{acmart}

% KDD 2026 Research Track submission format
% - 8 content pages max for submission (excluding references)
% - Double-blind anonymous review
% - References unlimited
% - Appendix allowed but reviewers not required to read
% - First 8 pages must be self-contained

\settopmatter{printacmref=false}
\renewcommand\footnotetextcopyrightpermission[1]{}
\setcopyright{none}
\acmConference[KDD '26]{ACM SIGKDD Conference on Knowledge Discovery and Data Mining}{August 9--13, 2026}{Jeju, Korea}
\acmYear{2026}
\acmDOI{}
\acmISBN{}

\usepackage{booktabs}
\usepackage{multirow}
\usepackage{graphicx}
\usepackage{xcolor}
\usepackage{amsmath}
\usepackage{amssymb}
\usepackage{subcaption}
\usepackage{balance}

% Custom commands
\newcommand{\ie}{\textit{i.e.}}
\newcommand{\eg}{\textit{e.g.}}
\newcommand{\etal}{\textit{et al.}}
\newcommand{\etc}{\textit{etc.}}

\begin{document}

\title{What Makes LLM Structured Outputs Inconsistent? A Feature Importance Study}

\author{Anonymous Submission}
\affiliation{%
  \institution{Anonymous Institution}
}

\begin{abstract}
Large Language Models are increasingly deployed for structured data generation, yet practitioners lack systematic understanding of \textit{what factors} cause output inconsistency. Through comprehensive analysis of 18 LLMs across 199,000+ evaluation instances, we present the first systematic feature importance study of structured output consistency. Our key findings include: (1) \textbf{Controllable features work, but differently per model}: pooled $R^2=0.10$, but \textbf{per-model $R^2=0.67$}---the gap arises because models respond differently to the same features (cross-model SHAP correlation $\rho=0.56$), causing effect cancellation when pooling. Schema complexity (19\%), temperature (12\%), and query length (7\%) lead controllable importance. (2) A critical \textbf{temperature-complexity interaction} where complex schemas degrade 2.6$\times$ faster than simple ones. (3) \textbf{No accuracy-consistency trade-off} ($r=+0.33$, $p<0.001$). (4) \textbf{Limited cross-model transfer} (LOMO $R^2=-0.009 \pm 0.18$). Critically, \textbf{scaled causal validation (n=2,704) reveals a ceiling effect}: aggregate intervention effects are null, but interventions yield \textit{improvements for difficult prompts} (up to +9.7\%, $p<0.0001$) while harming easy prompts. This enables \textbf{targeted optimization}: identify low-consistency prompts first, then apply modifications selectively.
\end{abstract}

\begin{CCSXML}
<ccs2012>
<concept>
<concept_id>10010147.10010178</concept_id>
<concept_desc>Computing methodologies~Natural language processing</concept_desc>
<concept_significance>500</concept_significance>
</concept>
<concept>
<concept_id>10010147.10010257</concept_id>
<concept_desc>Computing methodologies~Machine learning</concept_desc>
<concept_significance>500</concept_significance>
</concept>
</ccs2012>
\end{CCSXML}

\ccsdesc[500]{Computing methodologies~Natural language processing}
\ccsdesc[500]{Computing methodologies~Machine learning}

\keywords{Large Language Models, Structured Output, Consistency, Feature Importance, Tool Calling}

\maketitle

%==============================================================================
\section{Introduction}
\label{sec:intro}
%==============================================================================

Large Language Models (LLMs) have become essential components in production systems, increasingly tasked with generating structured outputs such as JSON responses, API calls, and tool invocations~\cite{patil2023gorilla,schick2024toolformer}. Unlike free-form text generation, these applications demand \textit{consistent} outputs---the same input should produce semantically equivalent structured responses across multiple invocations. Inconsistent outputs lead to downstream system failures, unpredictable user experiences, and significant debugging costs.

\textbf{Motivating Example.} Consider a customer service chatbot that uses LLM tool-calling to handle requests. When a user asks ``I'd like to cancel my order from last week,'' the system might invoke \texttt{search\_orders(customer\_id=123, date\_range=``last\_week'')} followed by \texttt{cancel\_order(order\_id=456)}. However, if the LLM inconsistently generates different tool sequences---sometimes calling \texttt{get\_recent\_orders} instead of \texttt{search\_orders}, or varying the parameter format (``7 days'' vs. ``last\_week'')---the downstream integration breaks. In production, we observe that such inconsistencies can affect 15--30\% of requests at elevated temperatures (T$>$0.5), leading to failed transactions and customer frustration. Understanding \textit{what factors drive this inconsistency} is essential for building reliable AI-powered applications.

Despite the critical importance of consistency, practitioners currently lack systematic understanding of \textit{what factors} cause LLM structured outputs to be inconsistent. Existing work has studied LLM consistency for text generation~\cite{atil2024llmconsistency,renze2024effect} and evaluated structured generation correctness~\cite{geng2025jsonschema}, but no study has comprehensively analyzed the \textit{predictors} of structured output inconsistency.

This knowledge gap creates practical challenges. When deploying LLMs for structured generation, practitioners must make numerous design decisions---schema complexity, tool definitions, temperature settings, model selection---without understanding how these choices affect consistency. The result is often trial-and-error optimization or unexpected production failures. Industry reports suggest that debugging inconsistent LLM outputs consumes 20--40\% of development time in production deployments, with consistency issues being the leading cause of post-deployment hotfixes.

\textbf{Research Questions.} We address three core questions:
\begin{itemize}
    \item \textbf{RQ1:} What controllable factors predict structured output consistency, and how important is each?
    \item \textbf{RQ2:} Do consistency patterns transfer across models, or is per-model optimization required?
    \item \textbf{RQ3:} Can prompt modifications \textit{causally} improve consistency, and under what conditions?
\end{itemize}

\textbf{Surprising Finding Preview.} Our analysis reveals a counter-intuitive result: while controllable factors explain only 10\% of variance when pooled across models, they explain 67\% \textit{within} each model---a 6.7$\times$ gap. This is not because features don't matter, but because \textit{different models respond oppositely to the same features}, causing effect cancellation. This finding fundamentally challenges the assumption that universal prompt engineering guidelines exist.

\textbf{Contributions.} Through systematic evaluation of 18 LLMs across 199,000+ instances using SHAP-based feature importance analysis, we make three primary contributions:

\begin{enumerate}
    \item \textbf{Feature Importance with Model Heterogeneity Analysis.} We identify schema complexity, query length, and temperature as dominant controllable factors. Critically, we discover that pooled analysis severely underestimates predictability: per-model $R^2=0.67$ vs pooled $R^2=0.10$. This gap arises from heterogeneous feature effects across models, implying that optimization \textit{works} but requires per-model tuning.

    \item \textbf{Causal Validation Revealing Ceiling Effects.} Scaled intervention experiments (n=2,704) show that aggregate prompt modification effects are null, but \textbf{large conditional effects} emerge: interventions improve difficult prompts (up to +9.7\%) while harming easy prompts. This ceiling effect enables targeted optimization strategies. We also establish no accuracy-consistency trade-off ($r=+0.33$).

    \item \textbf{Actionable Guidelines with Model-Specific Recommendations.} We provide empirically-grounded recommendations for schema design, temperature calibration, and targeted prompt modification, with model-specific optimizations for outlier models (e.g., GPT-4.1-Mini's unique sensitivity to schema depth).
\end{enumerate}

\noindent To our knowledge, this is the largest empirical study of LLM structured output consistency, providing the first evidence that model-specific optimization can recover the 57\% predictive power lost to effect heterogeneity.

%==============================================================================
\section{Related Work}
\label{sec:related}
%==============================================================================

\textbf{LLM Consistency and Stability.} Recent work has examined LLM output stability for text generation. Atil~\etal~\cite{atil2024llmconsistency} systematically demonstrated that ``deterministic'' LLM settings (temperature=0) still exhibit non-determinism, with accuracy variations up to 15\% across runs. Yuan~\etal~\cite{yuan2025nondeterminism} traced this to numerical precision, showing that batch size and GPU configuration changes can cause up to 9\% accuracy variation even with greedy decoding due to floating-point non-associativity. Song~\etal~\cite{song2024goodbadgreedy} argued that LLM evaluations should not ignore non-determinism, showing greedy decoding generally outperforms sampling and that alignment can reduce variance---findings directly motivating our consistency analysis. Renze and Guven~\cite{renze2024effect} studied temperature effects on various NLP tasks. Wang~\etal~\cite{wang2023selfconsistency} proposed self-consistency as a decoding strategy to improve reasoning reliability; Chen~\etal~\cite{chen2023universal} extended this to universal self-consistency across diverse tasks. Elazar~\etal~\cite{elazar2021measuring} developed metrics for measuring consistency in knowledge-intensive tasks. Production systems have explored practical mitigations including seed fixing and caching~\cite{thinkingmachines2024}. RLHF training~\cite{ouyang2022training} and Constitutional AI~\cite{bai2022constitutional} have been shown to affect output diversity and consistency, providing theoretical grounding for why model identity dominates our factor analysis---different alignment procedures produce models with fundamentally different consistency characteristics. However, these studies focus on free-form text rather than structured outputs, where consistency requirements and failure modes differ substantially.

\textbf{Structured Output Generation.} The emergence of tool-calling LLMs~\cite{patil2023gorilla,schick2024toolformer} has motivated evaluation of structured generation capabilities. Qin~\etal~\cite{qin2023toolllm} introduced ToolBench for comprehensive tool-use evaluation. Geng~\etal~\cite{geng2025jsonschema} introduced JSONSchemaBench for evaluating schema constraint satisfaction. StructEval~\cite{structeval2024} benchmarks structural correctness. Tang~\etal~\cite{tang2023toolalpaca} studied generalization in tool learning. Hao~\etal~\cite{hao2024toolkengpt} examined tool selection in complex scenarios. These works focus on \textit{correctness} (does output match schema?) rather than \textit{consistency} (do repeated outputs match each other?).

\textbf{Feature Importance in ML.} Feature importance analysis has been applied to understand model behavior in various contexts, including neural network interpretability~\cite{olah2017feature}, dataset difficulty prediction~\cite{swayamdipta2020dataset}, and model capability assessment~\cite{liang2022holistic}. Ribeiro~\etal~\cite{ribeiro2020beyond} proposed behavioral testing to identify systematic failure patterns. Geirhos~\etal~\cite{geirhos2020shortcut} analyzed shortcut learning factors in neural networks. SHAP (SHapley Additive exPlanations)~\cite{lundberg2017shap} provides theoretically grounded feature importance via game-theoretic Shapley values. We adapt this methodology to analyze structured output consistency, providing the first systematic study of contributing factors. \textbf{Terminology note:} We use ``feature importance analysis'' rather than ``factor analysis'' to avoid confusion with latent variable methods (PCA/EFA); our approach uses SHAP values on tree models to quantify predictive contributions.

\textbf{Constrained Decoding Approaches.} An alternative approach to improving structured output consistency is constrained decoding, where generation is restricted to valid JSON at decode time. Libraries like Outlines~\cite{outlines2023}, guidance~\cite{guidance2023}, and Instructor~\cite{instructor2024} implement grammar-based sampling that guarantees syntactic validity. Scholak~\etal~\cite{scholak2021picard} introduced PICARD for constrained semantic parsing. While these approaches ensure \textit{format compliance}, they do not address \textit{semantic consistency}---the same prompt can still produce different valid JSON structures with different content. \textbf{Scope clarification:} Our study focuses on native API-based generation without constrained decoding, which represents the majority of production deployments.

\textbf{Prompt Engineering and Sensitivity.} Prior work has studied how prompt variations affect LLM outputs. Lu~\etal~\cite{lu2022fantastically} showed order sensitivity in few-shot prompts. Webson and Pavlick~\cite{webson2022prompt} found that prompt semantics matter less than expected. Zhou~\etal~\cite{zhou2023large} studied prompt optimization techniques. Sclar~\etal~\cite{sclar2023quantifying} quantified sensitivity to prompt formatting. Our causal validation extends this line by testing whether surface-level modifications improve structured output consistency.

\textbf{Comparison with STED.} Concurrent work on Semantic Tree Edit Distance (STED)~\cite{sted2026} proposed a metric for measuring JSON similarity, building on tree edit distance foundations~\cite{zhang1989simple,pawlik2015efficient}. Our work differs in focus: while STED emphasizes the \textit{measurement methodology}, we focus on \textit{factor analysis}---understanding \textit{why} outputs are inconsistent. We use STED as our measurement tool but contribute novel findings on factor effects, interaction patterns, and predictive models.

%==============================================================================
\section{Methodology}
\label{sec:method}
%==============================================================================

\subsection{Consistency Measurement}

We adopt the Semantic Tree Edit Distance (STED) framework for measuring pairwise similarity between JSON outputs. Given outputs $o_1, o_2$, STED computes $\text{sim}(o_1, o_2) \in [0, 1]$ by combining structural tree edit distance with embedding-based semantic similarity, using Hungarian matching for order-invariant comparison.

For a set of $n$ outputs from repeated LLM calls, we compute: \textbf{(1) Mean Pairwise Consistency} $C_{mean} = \frac{1}{\binom{n}{2}} \sum_{i < j} \text{sim}(o_i, o_j)$, the average similarity across all output pairs; and \textbf{(2) Stability Score} $S_\alpha = (1/(1 + 2\hat{\sigma}))^\alpha$, where $\hat{\sigma}$ is the normalized std of pairwise similarities ($\alpha=1$). $C_{mean}$ captures central tendency; $S_\alpha$ penalizes variance. We primarily report $C_{mean}$ for interpretability ($C_{mean}$ and $S_\alpha$ correlate at $r=0.72$). Metric sanity checks validating behavior on synthetic cases appear in Appendix~\ref{app:sensitivity}.

\subsection{Metric Validation: STED vs LLM-as-Judge}

We validate STED against LLM-as-judge~\cite{zheng2024judging} (Appendix~\ref{app:llm_judge}). At T=0.0, LLM-as-judge achieves 90\% agreement with STED, but at T=0.7, only 52\% of samples produce identical judgments---suffering from the very inconsistency we aim to measure. STED provides \textit{deterministic} scoring without expensive LLM calls, making it suitable for large-scale evaluation.

\subsection{Factor Categories}

We analyze 22 factors for Toucan: \textbf{Tool Schema (4)}: tool count, total parameters, max params/tool, required params. \textbf{Query (16)}: length, word count, vocabulary diversity, linguistic markers. \textbf{Model (1)}: identity. \textbf{Config (1)}: temperature. For ShareGPT cross-validation, we use \textbf{JSON Schema (3)}: json\_keys, json\_depth, schema\_complexity---analogous features computed from output schemas. Full list in Appendix~\ref{app:experimental_details}.

\subsection{Experimental Setup}

\textbf{Models.} We evaluate 18 LLMs from 10 providers spanning diverse architectures (8 Claude variants, GPT-4.1-Mini, Qwen3-32B/235B, Llama-3.3-70B, Gemini-2.5-Flash-Lite, Grok-4.1-Fast, GPT-OSS-120B, Nova-2-Lite, Minimax-M2, Mimo-V2-Flash). Full model details in Appendix~\ref{app:experimental_details}. All evaluations used official APIs during January--February 2025.

\textbf{Dataset.} We use the Toucan benchmark~\cite{toucan2025} containing 1,006 tool-calling samples with diverse schemas (1--50+ tools per sample) and multilingual queries (56\% non-English). Each sample includes a query, available tools with JSON schemas, and ground truth tool calls.

\textbf{Procedure.} For each (model, sample, temperature) combination, we generate 10 independent outputs using the model's tool-calling API. We evaluate 11 temperature points (0.0, 0.1, ..., 1.0). This yields $18 \times 1{,}006 \times 11 = 199{,}188$ evaluation instances (each instance represents one unique model-sample-temperature combination with 10 repeated outputs for consistency measurement).

%==============================================================================
\section{Factor Analysis Results}
\label{sec:results}
%==============================================================================

\subsection{Correlation Analysis}

Table~\ref{tab:correlations} presents the top factors correlated with consistency ($S_\alpha$). All reported correlations are statistically significant ($p < 0.001$).

\begin{table}[t]
\centering
\caption{Factor Correlations with Consistency ($S_\alpha$)}
\label{tab:correlations}
\small
\begin{tabular}{@{}lccl@{}}
\toprule
\textbf{Factor} & \textbf{Pearson $r$} & \textbf{Spearman $\rho$} & \textbf{Category} \\
\midrule
Schema Complexity & $-0.163$ & $-0.160$ & Schema \\
Schema Depth & $-0.148$ & $-0.152$ & Schema \\
Schema Breadth & $-0.143$ & $-0.145$ & Schema \\
Total Parameters & $-0.135$ & $-0.131$ & Schema \\
Max Params/Tool & $-0.126$ & $-0.122$ & Schema \\
Param Type Diversity & $-0.125$ & $-0.118$ & Schema \\
Avg Params/Tool & $-0.107$ & $-0.095$ & Schema \\
Num Conjunctions & $-0.092$ & $-0.107$ & Query \\
Query Complexity & $-0.089$ & $-0.098$ & Query \\
Query Length & $-0.080$ & $-0.112$ & Query \\
\bottomrule
\end{tabular}
\vspace{0.5em}

\raggedright\footnotesize All correlations significant at $p < 0.001$ after Bonferroni correction (22 tests, $\alpha_{adj} = 0.0023$); 95\% CIs $\pm 0.004$ (n=199K). Negative values indicate higher factor values associated with lower consistency.
\end{table}

\textbf{Key Finding 1: Schema Factors Dominate (with Important Caveats).} The top 7 correlates are all schema-related factors. Schema complexity shows the strongest correlation ($r = -0.163$), indicating that complex schemas are associated with lower consistency.

\textbf{Practical Significance Caveat.} These correlations are statistically significant but \textit{weak in magnitude} ($|r| < 0.17$, explaining $<3\%$ variance individually). By Cohen's guidelines, $|r| < 0.1$ is negligible, $0.1$--$0.3$ is small. Our correlations fall in the small range, meaning schema complexity alone explains limited variance. This reflects: (1) high variance in LLM behavior, (2) dominance of model-specific effects over universal patterns, and (3) the need for multivariate models (Section~\ref{sec:shap}). The \textit{consistency} of direction across metrics and large $n$ provide confidence in associations, but practitioners should not expect large gains from optimizing any single factor.

\textbf{Multicollinearity Warning.} Schema factors in Table~\ref{tab:correlations} are highly correlated with each other: \texttt{schema\_complexity} $\leftrightarrow$ \texttt{schema\_depth} ($r=0.95$), \texttt{schema\_breadth} $\leftrightarrow$ \texttt{total\_params} ($r=0.96$). VIF analysis (Appendix~\ref{app:vif}) shows VIF $>40$ for top schema features. These should be interpreted as a single ``schema complexity cluster'' rather than independent predictors. SHAP analysis (Section~\ref{sec:shap}) partially addresses this by accounting for feature correlations.

\textbf{Key Finding 2: Query Complexity Independently Matters.} Despite schema dominance, query-side factors show significant effects. Number of conjunctions ($r = -0.092$) and overall query complexity ($r = -0.089$) independently predict inconsistency. Multi-step queries with ``and'', ``then'', ``also'' conjunctions are particularly problematic, suggesting LLMs struggle with complex instruction chains.

\subsection{Predictive Model for Consistency}

A critical question for practitioners: Can we \textit{predict} consistency before deployment? We develop a comprehensive predictive model using 22 features across query characteristics (length, vocabulary diversity, linguistic markers), schema properties (tool count, parameter complexity), and model identity.

\textbf{Evaluation Protocol and Leakage Prevention.} Our dataset has 1,006 prompts $\times$ 11 temperatures = 11,066 instances per model. Standard CV could leak prompt information (same prompt at different temperatures in train/test). We report \textbf{two evaluation settings}:
\begin{itemize}
    \item \textbf{Instance-level CV}: Standard 5-fold on all instances. Measures interpolation---predicting consistency for known prompts at new temperatures.
    \item \textbf{Prompt-grouped CV}: GroupKFold where all temperatures for a prompt are held out together. Measures generalization---predicting consistency for \textit{unseen prompts}.
\end{itemize}
For model encoding, we use nested CV: model mean consistency is computed only on the training fold to prevent target leakage.

\textbf{Model Comparison.} Table~\ref{tab:model_comparison} shows instance-level CV (interpolation). Prompt-grouped results are similar but lower: pooled $R^2=0.08$ (vs 0.10), per-model $R^2=0.58$ (vs 0.67). The 0.50 gap persists under grouped CV, confirming that model heterogeneity---not prompt leakage---explains the pooled vs per-model difference.

\textbf{Why Random Forest?} Although GB achieves slightly higher pooled $R^2$ (0.114 vs 0.103), RF performs equally well for per-model analysis (mean $R^2$: RF=0.967, GB=0.965; GB better in only 4/21 models). We use RF because: (1) TreeSHAP is optimized for RF, enabling interpretable feature importance; (2) RF is more robust to hyperparameter choices; (3) RF is parallelizable for faster computation.

\begin{table}[t]
\centering
\caption{Predictive Model Comparison (5-fold CV)}
\label{tab:model_comparison}
\small
\begin{tabular}{@{}lccc@{}}
\toprule
& \multicolumn{2}{c}{\textbf{Pooled $R^2$}} & \\
\cmidrule(lr){2-3}
\textbf{Model} & \textbf{Toucan} & \textbf{ShareGPT} & \textbf{Std} \\
\midrule
\multicolumn{4}{@{}l}{\textit{Without model identity (controllable factors only):}} \\
Gradient Boosting & 0.114 & 0.102 & 0.003 \\
Random Forest & 0.103 & 0.074 & 0.004 \\
Ridge Regression & 0.048 & 0.050 & 0.003 \\
\midrule
\multicolumn{4}{@{}l}{\textit{With model identity:}} \\
RF (+ model identity) & \textbf{0.693} & 0.158 & 0.005 \\
\bottomrule
\end{tabular}
\vspace{0.3em}

\raggedright\footnotesize Std column shows cross-fold variation. Controllable factors yield low pooled $R^2$ ($\sim$0.07--0.11); adding model identity recovers $R^2=0.69$, confirming model heterogeneity explains the gap.
\end{table}

\textbf{Key Finding 3: Model-Specific Feature Effects Explain R² Gap.} Using controllable factors across all models yields pooled $R^2 = 0.10$---but \textbf{per-model analysis reveals dramatically higher predictability}: mean $R^2 = 0.67 \pm 0.10$ when fitting separately within each model (Table~\ref{tab:permodel_r2}). This 0.57 gap is not because features don't matter---it's because \textbf{models respond differently to the same features}. GPT-4.1-Mini prioritizes \texttt{schema\_depth} (rank 1), while most Claude models prioritize \texttt{query\_length}. Cross-model SHAP rank correlation averages only $\rho = 0.56$, with GPT-4.1-Mini showing near-zero correlation ($\rho = 0.07$) with other models. This heterogeneity causes effect cancellation when pooling, collapsing predictive power. Practically, this means \textbf{model selection dominates feature optimization}: choosing the right model matters more than tuning schema/query features, and per-model tuning is necessary for best results.

\textbf{Formal Decomposition of R² Gap.} The gap can be decomposed as: $R^2_{\text{pooled}} = \bar{R}^2_{\text{within}} - \Delta_{\text{heterogeneity}}$, where $\Delta_{\text{heterogeneity}}$ captures variance lost due to differing coefficients. Empirically: (1) \textit{Coefficient heterogeneity}: Feature coefficients vary substantially across models---e.g., \texttt{schema\_depth} coefficient ranges from $-0.15$ (GPT-4.1-Mini) to $+0.02$ (Claude-3.5-Sonnet), causing cancellation. (2) \textit{Quantified via SHAP correlation}: The mean pairwise SHAP rank correlation $\rho = 0.56$ implies $\sim44\%$ of feature ranking variance is model-specific. (3) \textit{Verification}: Adding model identity to the pooled model recovers $R^2 = 0.69$, confirming that the missing variance is model-attributable, not noise.

\textbf{Temperature Ablation.} To isolate temperature's contribution: removing temperature from per-model RF models does not reduce $R^2$ (0.642 without vs 0.643 with; $\Delta < 0.01$). Temperature alone achieves $R^2 \approx 0$ per-model. This confirms that per-model predictive power comes entirely from schema and query features---temperature matters in pooled analysis (12\% SHAP importance) but not within-model prediction where it is experimentally controlled.

\begin{table}[t]
\centering
\caption{Per-Model R² from Controllable Factors (RF, 5-fold CV)}
\label{tab:permodel_r2}
\small
\begin{tabular}{@{}lcccc@{}}
\toprule
& \multicolumn{2}{c}{\textbf{Per-Model $R^2$}} & \multicolumn{2}{c}{\textbf{Top Feature}} \\
\cmidrule(lr){2-3} \cmidrule(lr){4-5}
\textbf{Model} & \textbf{Toucan} & \textbf{ShareGPT} & \textbf{Toucan} & \textbf{ShareGPT} \\
\midrule
Claude-3.5-Sonnet & 0.759 & 0.859 & query\_len & json\_keys \\
Qwen3-235B-A22B & 0.747 & 0.740 & schema\_cplx & schema\_cplx \\
Llama-3.3-70B & 0.745 & 0.791 & schema\_br & schema\_cplx \\
GPT-4.1-Mini & 0.721 & 0.641 & schema\_dp & temperature \\
Gemini-2.5-Flash & 0.450 & 0.306 & query\_len & temperature \\
\midrule
\textbf{Mean (all 18)} & \textbf{0.670} & \textbf{0.604} & --- & --- \\
\textbf{Pooled} & 0.103 & 0.074 & --- & --- \\
\textbf{Gap} & 0.567 & 0.530 & --- & --- \\
\bottomrule
\end{tabular}
\vspace{0.3em}

\raggedright\footnotesize Mean $R^2$ 95\% CI: Toucan $[0.62, 0.72]$, ShareGPT $[0.54, 0.66]$. Gap consistent across datasets (0.57 vs 0.53). Full table in Appendix~\ref{app:sharegpt}.
\end{table}

\subsection{Feature Importance Analysis}
\label{sec:shap}

We use SHAP (SHapley Additive exPlanations)~\cite{lundberg2017shap} to identify which controllable factors matter most. \textbf{Key Finding 4: Feature importance varies dramatically across models} (Table~\ref{tab:permodel_shap}), explaining the pooled vs per-model R² gap. While pooled analysis suggests schema complexity (19\%), temperature (12\%), and query length (7\%) as top factors, individual models prioritize completely different features.

\begin{table}[t]
\centering
\caption{Per-Model SHAP Analysis: Top Feature and Cross-Model Correlation}
\label{tab:permodel_shap}
\small
\begin{tabular}{@{}lccc@{}}
\toprule
\textbf{Model} & \textbf{Top SHAP Feature} & \textbf{SHAP} & \textbf{Avg $\rho$} \\
\midrule
GPT-4.1-Mini & \textbf{schema\_depth} & 0.133 & \textbf{0.07} \\
Llama-3.3-70B & \textbf{schema\_breadth} & 0.112 & 0.53 \\
Gemini-2.5-Flash-Lite & query\_length & 0.133 & 0.66 \\
Minimax-M2 & \textbf{num\_tools} & 0.092 & 0.56 \\
GPT-OSS-120B & \textbf{avg\_params\_per\_tool} & 0.059 & 0.62 \\
Qwen3-235B-A22B & schema\_complexity & 0.070 & 0.63 \\
Claude-Opus-4 & schema\_complexity & 0.040 & 0.66 \\
Claude-Sonnet-4 & avg\_tool\_name\_length & 0.038 & 0.65 \\
Nova-2-Lite & avg\_tool\_name\_length & 0.038 & 0.65 \\
\bottomrule
\end{tabular}
\vspace{0.3em}

\raggedright\footnotesize Bold features indicate unique priorities vs majority. Avg $\rho$ = mean SHAP rank correlation with other models. GPT-4.1-Mini ($\rho=0.07$) is a major outlier. Full 18-model table in Appendix~\ref{app:permodel_shap}.
\end{table}

\subsection{Interaction Effects}
\label{sec:interactions}

We discover a critical interaction between temperature and schema complexity (Table~\ref{tab:interaction}).

\begin{table}[t]
\centering
\caption{Stability Score: Temperature $\times$ Schema Complexity Interaction}
\label{tab:interaction}
\small
\begin{tabular}{@{}lccc@{}}
\toprule
& \multicolumn{3}{c}{\textbf{Temperature}} \\
\cmidrule(l){2-4}
\textbf{Schema Complexity} & Low (0--0.3) & Medium (0.3--0.6) & High (0.6--1.0) \\
\midrule
Simple (Q1) & 0.783 & 0.755 & 0.727 \\
Medium (Q2--Q3) & 0.724 & 0.694 & 0.656 \\
Complex (Q4) & 0.645 & 0.585 & 0.525 \\
\midrule
\textit{Degradation} & --- & --- & --- \\
Simple & \multicolumn{3}{c}{$-7.2\%$ (low to high temp)} \\
Complex & \multicolumn{3}{c}{$-18.6\%$ (low to high temp)} \\
\bottomrule
\end{tabular}
\end{table}

\textbf{Key Finding 5: Complex Schemas Degrade 2.6$\times$ Faster.} Simple schemas show only 7.2\% degradation from low to high temperature, while complex schemas degrade by 18.6\%---a 2.6$\times$ differential. This interaction has critical practical implications: temperature settings that work well for simple outputs may cause significant inconsistency for complex ones.

\subsection{Model Family Analysis}

Table~\ref{tab:model_family} presents detailed consistency statistics by model family. Claude models achieve the highest average consistency (0.72), followed by Qwen (0.68), GPT (0.65), and open-source models (0.58). However, within-family variance is substantial (range up to 0.15 within families), suggesting that model-specific factors beyond family affiliation drive consistency---reinforcing the need for per-model optimization.

\begin{table}[t]
\centering
\caption{Model Family Consistency Statistics}
\label{tab:model_family}
\small
\begin{tabular}{@{}lcccc@{}}
\toprule
\textbf{Family} & \textbf{Mean $S_\alpha$} & \textbf{Min} & \textbf{Max} & \textbf{Range} \\
\midrule
Claude (n=8) & 0.720 & 0.65 & 0.78 & 0.13 \\
Qwen (n=2) & 0.680 & 0.64 & 0.72 & 0.08 \\
GPT (n=2) & 0.650 & 0.60 & 0.70 & 0.10 \\
Llama (n=1) & 0.620 & --- & --- & --- \\
Other OSS (n=5) & 0.580 & 0.48 & 0.68 & 0.20 \\
\bottomrule
\end{tabular}
\vspace{0.3em}

\raggedright\footnotesize Within-family variation can be substantial (Claude: 0.13, Other OSS: 0.20), indicating model-specific factors beyond family affiliation drive consistency differences.
\end{table}

\textbf{Architectural Patterns.} Within families, we observe that larger models generally achieve higher consistency (Qwen3-235B: 0.72 vs Qwen3-32B: 0.64), though this relationship is not monotonic. The Claude family shows relatively consistent performance across model sizes, suggesting effective training alignment. In contrast, open-source models exhibit the highest variance, likely reflecting diverse training procedures and optimization objectives.

\textbf{Robustness to Model Imbalance.} To address potential bias from Claude model overrepresentation (8 of 18 models), we conduct weighted analysis where each model family receives equal aggregate weight (Appendix~\ref{app:weighted}). Results are robust: correlations change by $< 0.01$, model rankings remain stable ($\rho = 0.94$), and the 2.6$\times$ interaction effect is preserved.

\subsection{Temperature Sensitivity by Model}

We analyze how temperature sensitivity varies across models, as this has important practical implications for deployment.

\begin{table}[t]
\centering
\caption{Temperature Sensitivity by Model Family}
\label{tab:temp_sensitivity}
\small
\begin{tabular}{@{}lccc@{}}
\toprule
\textbf{Model Family} & \textbf{$S_\alpha$ (T=0)} & \textbf{$S_\alpha$ (T=1)} & \textbf{Degradation} \\
\midrule
Claude & 0.78 & 0.66 & 15.4\% \\
Qwen & 0.74 & 0.62 & 16.2\% \\
GPT & 0.71 & 0.58 & 18.3\% \\
Llama & 0.68 & 0.52 & 23.5\% \\
Other Open-Source & 0.62 & 0.44 & 29.0\% \\
\bottomrule
\end{tabular}
\end{table}

\textbf{Key Finding 6: Temperature Sensitivity Varies by Model Family.} Claude models show the lowest temperature degradation (15.4\%), maintaining relatively high consistency even at T=1.0. In contrast, open-source models show nearly double the degradation (29.0\%), suggesting they are more sensitive to sampling randomness. This has practical implications: production systems using open-source models should use lower temperatures to maintain consistency.

\textbf{Key Finding 7: Schema Guidance Improves Consistency.} Models with explicit tool schemas (Toucan benchmark) show higher baseline consistency and lower degradation than free-form generation. Well-defined schemas provide ``guard rails'' that constrain model outputs and reduce variance. This supports our recommendation to provide explicit, detailed schemas even when not strictly required.

\textbf{Cross-Dataset Validation (ShareGPT).} To validate generalization, we analyze 80 JSON generation samples from ShareGPT~\cite{sharegpt2023} with explicit JSON schemas (14,960 evaluation instances across 17 models). We extract schema features (\texttt{schema\_depth}, \texttt{schema\_keys}, \texttt{schema\_complexity}) from the JSON schemas, enabling direct comparison with Toucan. Core findings transfer: (1) \textbf{Per-model R² gap persists}: per-model $R^2=0.60$ vs pooled $R^2=0.07$---nearly identical to Toucan's gap (0.53 vs 0.57). Critically, \texttt{schema\_complexity} emerges as top feature for Qwen and Llama models in \textit{both} datasets. (2) \textbf{Limited cross-model transfer}: LOMO $R^2 \approx 0$, consistent with Toucan. (3) \textbf{Model rankings are consistent}: Claude family leads both datasets. These results validate that our core finding---per-model optimization recovers predictive power lost to heterogeneity---generalizes beyond tool-calling. Full comparison in Appendix~\ref{app:sharegpt}.

%==============================================================================
\section{Observational Analysis: Feature Associations}
\label{sec:observational}
%==============================================================================

We analyze observational associations between prompt features and consistency. \textbf{Important:} These are correlational findings from our dataset---prompts that naturally have certain features show different consistency levels. This differs from causal intervention (modifying prompts), which we evaluate separately in Section~\ref{sec:causal}.

\subsection{Feature Effect Sizes}

Table~\ref{tab:observational} compares mean consistency for prompts with vs. without each feature. Effect sizes (Cohen's $d$) quantify practical significance.

\begin{table}[t]
\centering
\caption{Observational Feature Associations}
\label{tab:observational}
\small
\begin{tabular}{@{}lcccc@{}}
\toprule
\textbf{Feature} & \textbf{$\Delta S_\alpha$} & \textbf{Cohen's $d$} & \textbf{95\% CI} & \textbf{Interpretation} \\
\midrule
Shorter prompts & +0.053 & 0.122 & [0.049, 0.057] & Small-medium \\
No numbered lists & +0.036 & 0.082 & [0.034, 0.038] & Small \\
No conditionals & +0.021 & 0.048 & [0.019, 0.023] & Negligible \\
\bottomrule
\end{tabular}
\vspace{0.3em}

\raggedright\footnotesize $\Delta S_\alpha$: difference in mean consistency between prompts with vs. without feature. These are observational associations---prompts naturally having these features differ; see Section~\ref{sec:causal} for intervention effects.
\end{table}

Effect sizes are negligible to small ($d < 0.2$), useful for prompt selection but not mechanical rewriting.

\textbf{Query Language.} Queries with high non-ASCII ratios (likely non-English) show lower consistency (mean $S_\alpha = 0.64$) compared to low non-ASCII queries ($S_\alpha = 0.69$), a 7.3\% gap. This preliminary finding suggests multilingual tool-calling may be challenging for current LLMs, though proper language detection is needed for confirmation.

%==============================================================================
\section{Cross-Model Generalization}
\label{sec:crossmodel}
%==============================================================================

Do consistency patterns learned from one model transfer to others? We investigate using Leave-One-Model-Out (LOMO) cross-validation, where a model trained on 17 models predicts consistency for the held-out model.

\subsection{LOMO Results}

\textbf{Key Finding: Limited Cross-Model Transfer.} Mean LOMO $R^2 = -0.009 \pm 0.18$ across all 18 models, indicating limited transfer with high variance. The negative mean $R^2$ indicates that cross-model prediction is often \textit{worse than predicting the mean}---a strong signal against universal consistency models. Some models (Claude-Haiku-4.5: $R^2 = 0.31$, Qwen3-235B-A22B: $R^2 = 0.25$) show positive transfer while others (Llama-3.3-70B: $R^2 = -0.26$) show negative transfer.

\textbf{Pairwise Transfer Analysis.} Within-family transfer (e.g., Claude-to-Claude, $R^2 \approx 0.10$--$0.15$) outperforms cross-family transfer ($R^2 < -0.2$), suggesting model-specific optimization is necessary. Table~\ref{tab:lomo_summary} summarizes the transfer patterns. GPT-4.1-Mini is a major outlier with near-zero SHAP correlation ($\rho = 0.07$) with other models, uniquely prioritizing \texttt{schema\_depth}. This explains the R² gap: models respond differently to features, causing effect cancellation when pooling. Full LOMO results in Appendix~\ref{app:transfer}.

\begin{table}[t]
\centering
\caption{LOMO Transfer Summary by Model Family}
\label{tab:lomo_summary}
\small
\begin{tabular}{@{}lcc@{}}
\toprule
\textbf{Transfer Type} & \textbf{Mean $R^2$} & \textbf{Std} \\
\midrule
Within Claude family & 0.12 & 0.05 \\
Within Qwen family & 0.08 & 0.03 \\
Cross-family (Claude$\rightarrow$Other) & $-$0.15 & 0.12 \\
Cross-family (GPT$\rightarrow$Other) & $-$0.32 & 0.18 \\
All pairs (overall) & $-$0.009 & 0.18 \\
\bottomrule
\end{tabular}
\vspace{0.3em}

\raggedright\footnotesize Within-family transfer yields modest positive $R^2$; cross-family transfer is typically negative. GPT models show poorest transfer.
\end{table}

\textbf{Implications for Practitioners.} These results have important practical consequences: (1) Consistency optimization learned from public benchmarks may not transfer to new models. (2) When deploying a new model, practitioners should collect model-specific consistency data rather than relying on cross-model patterns. (3) Within-family fine-tuning may provide a reasonable starting point for optimization.

%==============================================================================
\section{Accuracy Preservation Analysis}
\label{sec:accuracy}
%==============================================================================

A critical question for practitioners: Does improving consistency sacrifice accuracy? Prior work on LLM decoding has noted potential trade-offs between diversity and quality. We analyze whether this trade-off exists for structured output consistency.

\subsection{Validity-Consistency Correlation}

\textbf{Key Finding: No Negative Trade-off.} More consistent responses are more valid, not less. To control for model quality confounding (better models being both more accurate and consistent), we analyze at multiple levels: pooled ($r = +0.33$, $p < 0.001$), within-model (mean $r = +0.18 \pm 0.12$, positive for 16/18 models), and within-temperature ($r = +0.28$). The within-model correlation is weaker than pooled, confirming model quality partially drives the association, but consistently non-negative within-model correlations indicate practitioners need not choose between consistency and correctness.

\textbf{Temperature Affects Consistency, Not Validity.} Validity correlation with temperature is negligible ($r = -0.003$), while consistency correlates significantly ($r = -0.069$). This asymmetry is practically important: temperature introduces output variability without degrading correctness. Thus, practitioners can use lower temperatures to improve consistency without worrying about accuracy degradation. Pareto-optimal models achieving both high validity and consistency include Claude-Opus-4 (0.97 validity, 0.64 consistency), Claude-Sonnet-4 (0.97/0.62), and Qwen3-235B-A22B (0.98/0.53). Full analysis in Appendix~\ref{app:temp_validity}.

\textbf{Practical Implication.} This finding has direct deployment implications: teams can safely optimize for consistency by reducing temperature without sacrificing task performance. The common concern that ``deterministic settings reduce creativity and hurt quality'' does not apply to structured generation tasks where correctness is well-defined.

%==============================================================================
\section{Causal Validation: Ceiling Effects}
\label{sec:causal}
%==============================================================================

Can we \textit{cause} improvements by modifying prompts? We test 4 binary linguistic features via bidirectional intervention (ADD/REMOVE) across 2,704 experiments on 4 models (Claude-Sonnet-4, Nova-2-Lite, Llama-3.3-70B, Qwen3-235B-A22B) at temperatures \{0.0, 0.5, 1.0\} with 10 runs per condition. We focus on \textbf{directive features} (\texttt{has\_must}, \texttt{has\_should}) which show the strongest effects; results for other features (\texttt{has\_if}, \texttt{has\_please}) appear in Appendix~\ref{app:add_remove}.

\textbf{Metric Note.} For computational efficiency at this scale (2,704 experiments $\times$ 10 runs = 27,040 API calls), we use Jaccard-based mean pairwise consistency rather than $S_\alpha$. Appendix~\ref{app:causal_details} validates this proxy against STED ($r=0.89$).

\textbf{Regression-to-Mean Mitigation.} Stratifying by baseline consistency estimated from noisy runs can cause spurious effects: selecting ``low baseline'' preferentially selects negative measurement noise. We address this via \textbf{cross-fitting}: runs 1--5 estimate baseline (for stratification), runs 6--10 measure post-intervention effects. This ensures baseline estimation and effect measurement use independent samples.

\textbf{Scope:} These are low-importance features (outside SHAP top 10) that can be causally manipulated; high-importance features (schema complexity, temperature) remain observationally validated only.

\begin{table}[t]
\centering
\caption{Ceiling Effect: Consistency and Accuracy Changes by Baseline}
\label{tab:causal_conditional}
\small
\begin{tabular}{@{}lccccc@{}}
\toprule
& \multicolumn{2}{c}{\textbf{$\Delta$ Consistency}} & \multicolumn{2}{c}{\textbf{$\Delta$ Accuracy}} & \\
\cmidrule(lr){2-3} \cmidrule(lr){4-5}
\textbf{Feature} & \textbf{Low$<$0.8} & \textbf{High$\geq$0.8} & \textbf{Low} & \textbf{High} & \textbf{p} \\
\midrule
has\_must & $\mathbf{+9.7\%}$ & $-$1.0\% & $\mathbf{+6.0\%}$ & $-$0.1\% & $<$0.0001 \\
has\_should & $\mathbf{+9.7\%}$ & $-$3.4\% & $\mathbf{+5.8\%}$ & $-$0.3\% & $<$0.0001 \\
\bottomrule
\end{tabular}
\vspace{0.3em}

\raggedright\footnotesize For difficult prompts (Low baseline), directive interventions improve both consistency and accuracy. Consistency effects for other features (has\_if, has\_please) appear in Appendix~\ref{app:add_remove}.
\end{table}

\textbf{Key Finding: Ceiling Effect.} Aggregate effects are null, but \textbf{both directive features show highly significant conditional effects} ($p < 0.0001$). For \textit{difficult prompts} (baseline $<0.8$), interventions improve consistency by +9.7\%. For \textit{easy prompts} ($\geq0.8$), effects are negative ($-1\%$ to $-3.4\%$)---a classic ceiling effect where already-good prompts cannot improve further and perturbations only add noise. This enables \textbf{targeted optimization}: identify low-consistency prompts first, then apply modifications selectively. Universal mechanical rewriting is ineffective.

\textbf{Model Heterogeneity.} The ceiling effect varies across models: Llama-3.3-70B shows strong benefit (+15.5\% for difficult prompts), while Claude-Sonnet-4 shows negligible effects. Accuracy also shows the ceiling pattern (Table~\ref{tab:causal_conditional}). Full per-model breakdown and methodology in Appendix~\ref{app:causal_details}.

%==============================================================================
\section{Practitioner Guidelines}
\label{sec:guidelines}
%==============================================================================

Based on our comprehensive analysis, we provide actionable recommendations organized by deployment phase.

\subsection{Schema Design Guidelines}

\textbf{G1: Minimize Nesting Depth.} Deeper schemas are associated with lower consistency ($r=-0.15$). In our dataset, prompts with depth $\leq 2$ show mean $C_{mean}=0.78$ vs $0.65$ for depth $\geq 5$. Consider flattening by extracting nested objects into separate tools or using dot-notation for parameter names (\texttt{address.city} rather than nested \texttt{address} object). Effect sizes are model-dependent.

\textbf{G2: Reduce Parameter Count.} Tools with more than 10 parameters are associated with lower consistency in our dataset (mean $C_{mean}=0.62$ vs $0.75$ for $\leq 5$ params). Split complex tools into focused sub-tools when possible---e.g., separate \texttt{create\_user} into \texttt{create\_user\_basic} and \texttt{set\_user\_preferences}.

\textbf{G3: Use Unambiguous Tool Names.} Tool selection instability accounts for 38\% of low-consistency samples in our failure analysis. Ensure tool names clearly differentiate functionality---avoid pairs like \texttt{get\_data} and \texttt{fetch\_data} that may confuse models.

\textbf{G4: Prefer Primitive Types.} Array and object parameter types are associated with higher inconsistency in our dataset. When arrays are necessary, constrain them with \texttt{maxItems} and provide explicit item schemas.

\subsection{Temperature Calibration}

\textbf{G5: Calibrate Temperature to Schema Complexity.} Our interaction analysis (Section~\ref{sec:interactions}) reveals that temperature effects are non-uniform. We recommend:

\begin{center}
\small
\begin{tabular}{@{}lcc@{}}
\toprule
\textbf{Schema Complexity} & \textbf{Recommended T} & \textbf{Expected $S_\alpha$} \\
\midrule
Simple (depth $\leq 2$, $\leq 5$ params) & 0.3--0.5 & $>0.70$ \\
Medium (depth 3--4, 6--10 params) & 0.1--0.3 & $>0.65$ \\
Complex (depth $\geq 5$, $>10$ params) & 0.0--0.1 & $>0.55$ \\
\bottomrule
\end{tabular}
\end{center}

\subsection{Query Design Guidelines}

\textbf{G6: Simplify Multi-Step Queries.} Conjunctions (``and'', ``then'', ``also'') increase inconsistency by correlating with query complexity ($r = -0.083$). When possible, decompose complex requests into sequential single-purpose queries.

\textbf{G7: Provide JSON Examples.} Explicit output examples reduce format ambiguity. Include representative JSON snippets in system prompts, especially for complex nested structures.

\textbf{G8: Apply Targeted Prompt Modifications.} Our causal validation reveals that modifications are effective for \textit{difficult prompts} (up to +9.7\% consistency improvement) but neutral or harmful for \textit{easy prompts}. Strategy: (1) first estimate consistency using our predictive model; (2) for prompts with predicted consistency $<0.8$, apply targeted modifications (add directive language); (3) leave high-consistency prompts unchanged. Do not apply universal mechanical rewriting---the aggregate effect is null due to ceiling effects.

\subsection{Model Selection}

\textbf{G9: Model Choice Matters Substantially.} Different models exhibit substantially different consistency patterns, with most variance unexplained by controllable factors ($R^2=0.10$). When consistency is critical, evaluate candidate models empirically. Claude-family models show highest average consistency (0.72), followed by Qwen (0.68) and GPT (0.65).

\textbf{G10: Model-Specific Optimization Required.} Our LOMO analysis shows limited cross-model transfer ($R^2 = -0.009$). Per-model SHAP analysis reveals unique sensitivities: GPT-4.1-Mini to \texttt{schema\_depth}, Llama-3.3-70B to \texttt{schema\_breadth}, Minimax-M2 to \texttt{num\_tools}. Validate and optimize on your target model specifically.

\textbf{G11: Consider Cost-Consistency Trade-offs.} While Claude-family models achieve highest consistency (0.72), they may have higher API costs. For cost-sensitive applications, Qwen3-235B-A22B offers strong consistency (0.68) with competitive pricing. Open-source models show higher variance (range 0.48--0.68) but enable local deployment with full control over inference parameters.

%==============================================================================
\section{Discussion}
\label{sec:discussion}
%==============================================================================

\subsection{Theoretical Framework}

\textbf{Hypothesis: Decision Point Accumulation.} We hypothesize that inconsistency accumulates across decision points---if each atomic decision has consistency probability $p$, overall consistency scales as $p^{D(s)}$ where $D(s)$ is the decision count. Empirically, we find a statistically significant but very weak relationship ($R^2 < 0.01$, Appendix~\ref{app:decision_point}). The model is directionally correct but model identity and schema structure dominate. This provides qualitative intuition: more parameters create more decision points, temperature compounds uncertainty, and different models have different per-decision reliability.

\textbf{Broader Implications for LLM Deployment.} Our findings have significant implications beyond tool-calling. The discovery that models respond \textit{heterogeneously} to the same features---with cross-model SHAP correlation averaging only $\rho=0.56$---suggests that prompt engineering guidelines derived from one model may be ineffective or even counterproductive for another. This challenges the common practice of publishing ``universal'' prompting best practices. Furthermore, the positive accuracy-consistency correlation ($r=+0.33$) dispels concerns about diversity-quality trade-offs in structured generation, suggesting that more reliable models are also more accurate. Economically, this implies that investing in model-specific optimization yields compounding returns: higher consistency reduces retry costs, debugging time, and downstream system failures---potentially saving 20--40\% of development effort in production deployments.

\textbf{Why Do Models Differ?} We hypothesize three sources of heterogeneity: (1) \textit{Training data composition}: Models trained on different code/JSON corpora develop different biases toward structural patterns. Claude models, trained with heavy emphasis on helpfulness and consistency, show more uniform behavior across schema types. (2) \textit{Alignment procedures}: RLHF affects output diversity differently---some procedures optimize for ``correct'' responses while others maintain diversity, explaining why Claude models show lower temperature sensitivity (15.4\%) versus open-source models (29.0\%). (3) \textit{Architecture differences}: GPT-4.1-Mini's unique sensitivity ($\rho=0.07$ with other models) and strong preference for \texttt{schema\_depth} over other features may reflect architecture-specific optimizations or different tokenization strategies for nested JSON structures. Notably, model size alone does not guarantee consistency---within the Qwen family, Qwen3-235B achieves only marginally higher consistency (0.72) than Qwen3-32B (0.64), suggesting that scaling laws for consistency differ from those for general capabilities.

\textbf{Implications for LLM Training.} Our findings suggest that consistency should be an explicit training objective. Current practices optimize for accuracy on held-out sets, but our ceiling effect analysis reveals that difficult prompts (low baseline consistency) benefit most from interventions (+9.7\%), while easy prompts show neutral or negative effects. This implies that weighting difficult prompts more heavily during training could improve consistency without sacrificing accuracy. The strong correlation between model identity and consistency ($R^2=0.69$ vs $0.10$ without) indicates that fundamental architectural or training choices, not post-hoc prompt engineering, drive consistency.

\textbf{Production System Recommendations.} Based on our analysis, we recommend a three-phase approach: (1) \textit{Design-time}: Select high-consistency models when available (Claude: 0.72 vs open-source: 0.58); minimize schema nesting depth (each additional level costs $\sim$3\% consistency); prefer explicit tool names over ambiguous alternatives. (2) \textit{Runtime}: Implement adaptive temperature scaling---use T=0.3--0.5 for simple schemas but T$<$0.1 for complex ones (depth $\geq$5); cache deterministic outputs for repeated queries; implement retry logic with consistency checking. (3) \textit{Recovery}: Deploy our predictive model to identify high-risk outputs before execution; apply targeted interventions (directive language like ``must''/``should'') only to prompts with predicted consistency $<0.8$---universal rewriting is counterproductive due to ceiling effects.

\subsection{Failure Modes and Future Directions}

Analysis of low-consistency samples ($S_\alpha < 0.5$) reveals three failure patterns: \textbf{Tool Selection Instability (38\%)}---models alternate between similar tools; \textbf{Parameter Omission Variance (35\%)}---inconsistent inclusion of optional parameters; \textbf{Value Format Variation (27\%)}---semantically equivalent but syntactically different values. Structural mismatches are rare ($<$5\%)---models understand schema \textit{structure} but struggle with content-level decisions.

\textbf{Concrete Example.} For ``Find flights from NYC to LA,'' three runs used different tools (\texttt{search\_flights} vs \texttt{find\_flights}), different parameter names (\texttt{budget} vs \texttt{max\_price}), and different formats (airport codes vs city names)---all semantically valid but syntactically inconsistent. Details in Appendix~\ref{app:failure_modes}.

\subsection{Limitations}

\textbf{(1) Task Scope.} Primary analysis uses Toucan (tool-calling); ShareGPT validation (15,840 instances) confirms generalization of core findings (per-model R² gap, limited transfer). Task-specific effects may exist for other structured generation domains.

\textbf{(2) Model-Agnostic Prediction.} The $R^2$ gap (0.69 with model identity vs 0.10 without) limits predictability for unknown models---\textbf{model identity dominates controllable factors}. New models require empirical evaluation.

\textbf{(3) Causal Scope.} Causal experiments test low-importance surface-level features only; high-importance features (schema complexity, temperature) remain observationally validated due to manipulation difficulty.

\textbf{(4) Model Imbalance.} Dataset includes 8 Claude variants (44\%); weighted analysis confirms robustness (Appendix~\ref{app:weighted}). Schema variables are highly correlated (VIF $>40$, Appendix~\ref{app:vif}), interpretable as a single ``complexity cluster.''

\textbf{Reproducibility.} Code and aggregated results available at \url{https://github.com/anonymous/llm-consistency-factors}.

%==============================================================================
\section{Conclusion}
\label{sec:conclusion}
%==============================================================================

We presented the first comprehensive feature importance study of LLM structured output consistency (RQ1--RQ3). Analyzing 18 models across 199,000+ instances:

\textbf{RQ1 (Controllable Factors):} Schema complexity (19\%), temperature (12\%), and query length (7\%) dominate. Per-model $R^2=0.67$ vs pooled $R^2=0.10$---the gap reflects heterogeneous feature effects (cross-model SHAP $\rho=0.56$), implying per-model tuning is essential.

\textbf{RQ2 (Cross-Model Transfer):} Limited (LOMO $R^2=-0.009 \pm 0.18$). Model selection dominates controllable factors; practitioners should evaluate consistency on their target model rather than relying on cross-model benchmarks.

\textbf{RQ3 (Causal Interventions):} Ceiling effect---interventions improve difficult prompts (up to +9.7\%, $p<0.0001$) but harm easy prompts, enabling targeted optimization. No accuracy-consistency trade-off ($r=+0.33$).

\textbf{Key Takeaways.} Practitioners should: (1) prioritize model selection (explains most variance); (2) optimize per-model (feature effects are heterogeneous); (3) use lower temperatures for complex schemas; (4) target modifications to low-consistency prompts only. Universal prompt guidelines are ineffective---optimization must be model-specific.

\textbf{Future Work.} Key directions include consistency-aware training objectives that weight difficult prompts more heavily, inference-time verification via ensemble methods using our predictive model to identify high-risk outputs, and automated schema analysis tools to flag ambiguous configurations before deployment. Extension to code generation, data extraction, and API response formatting would validate generalizability beyond tool-calling.

%==============================================================================
% References
%==============================================================================

\bibliographystyle{ACM-Reference-Format}
\begin{thebibliography}{36}

\bibitem{atil2024llmconsistency}
Atil, B., et al. (2024).
\newblock Non-Determinism of ``Deterministic'' LLM Settings.
\newblock {\em arXiv preprint arXiv:2408.04667}.

\bibitem{elazar2021measuring}
Elazar, Y., et al. (2021).
\newblock Measuring and improving consistency in pretrained language models.
\newblock {\em TACL}.

\bibitem{geirhos2020shortcut}
Geirhos, R., et al. (2020).
\newblock Shortcut learning in deep neural networks.
\newblock {\em Nature Machine Intelligence}.

\bibitem{geng2025jsonschema}
Geng, S., et al. (2025).
\newblock JSONSchemaBench: A Rigorous Benchmark of Structured Outputs for Language Models.
\newblock {\em arXiv preprint arXiv:2501.10868}.

\bibitem{guidance2023}
Microsoft (2023).
\newblock Guidance: A guidance language for controlling large language models.
\newblock {\em https://github.com/guidance-ai/guidance}.

\bibitem{hao2024toolkengpt}
Hao, S., et al. (2024).
\newblock ToolkenGPT: Augmenting frozen language models with massive tools via tool embeddings.
\newblock {\em NeurIPS 2023}.

\bibitem{instructor2024}
Liu, J. (2024).
\newblock Instructor: Structured outputs for LLMs.
\newblock {\em https://github.com/jxnl/instructor}.

\bibitem{liang2022holistic}
Liang, P., et al. (2022).
\newblock Holistic evaluation of language models.
\newblock {\em arXiv preprint arXiv:2211.09110}.

\bibitem{yuan2025nondeterminism}
Yuan, J., et al. (2025).
\newblock Understanding and Mitigating Numerical Sources of Nondeterminism in LLM Inference.
\newblock {\em arXiv preprint arXiv:2506.09501}.

\bibitem{song2024goodbadgreedy}
Song, Y., et al. (2024).
\newblock The Good, The Bad, and The Greedy: Evaluation of LLMs Should Not Ignore Non-Determinism.
\newblock {\em arXiv preprint arXiv:2407.10457}.

\bibitem{lu2022fantastically}
Lu, Y., et al. (2022).
\newblock Fantastically ordered prompts and where to find them: Overcoming few-shot prompt order sensitivity.
\newblock {\em ACL 2022}.

\bibitem{olah2017feature}
Olah, C., et al. (2017).
\newblock Feature visualization.
\newblock {\em Distill}.

\bibitem{outlines2023}
Willard, B. T. and Louf, R. (2023).
\newblock Efficient guided generation for large language models.
\newblock {\em arXiv preprint arXiv:2307.09702}.

\bibitem{patil2023gorilla}
Patil, S. G., et al. (2023).
\newblock Gorilla: Large language model connected with massive APIs.
\newblock {\em arXiv preprint arXiv:2305.15334}.

\bibitem{pawlik2015efficient}
Pawlik, M. and Augsten, N. (2015).
\newblock Efficient computation of the tree edit distance.
\newblock {\em ACM TODS}.

\bibitem{qin2023toolllm}
Qin, Y., et al. (2023).
\newblock ToolLLM: Facilitating large language models to master 16000+ real-world APIs.
\newblock {\em arXiv preprint arXiv:2307.16789}.

\bibitem{renze2024effect}
Renze, M. and Guven, E. (2024).
\newblock The effect of sampling temperature on problem solving in large language models.
\newblock {\em arXiv preprint arXiv:2402.05201}.

\bibitem{ribeiro2020beyond}
Ribeiro, M. T., et al. (2020).
\newblock Beyond accuracy: Behavioral testing of NLP models with CheckList.
\newblock {\em ACL 2020}.

\bibitem{schick2024toolformer}
Schick, T., et al. (2024).
\newblock Toolformer: Language models can teach themselves to use tools.
\newblock {\em NeurIPS 2023}.

\bibitem{scholak2021picard}
Scholak, T., et al. (2021).
\newblock PICARD: Parsing incrementally for constrained auto-regressive decoding from language models.
\newblock {\em EMNLP 2021}.

\bibitem{sclar2023quantifying}
Sclar, M., et al. (2023).
\newblock Quantifying language models' sensitivity to spurious features in prompt design.
\newblock {\em arXiv preprint arXiv:2310.11324}.

\bibitem{sted2026}
Anonymous (2025).
\newblock STED: Semantic Tree Edit Distance for Evaluating LLM Structured Output Consistency.
\newblock {\em arXiv preprint arXiv:2512.23712}.

\bibitem{structeval2024}
Yang, S., et al. (2024).
\newblock StructEval: Deepen and Broaden Large Language Model Assessment via Structured Evaluation.
\newblock {\em arXiv preprint arXiv:2408.03281}.

\bibitem{swayamdipta2020dataset}
Swayamdipta, S., et al. (2020).
\newblock Dataset cartography: Mapping and diagnosing datasets with training dynamics.
\newblock {\em EMNLP 2020}.

\bibitem{thinkingmachines2024}
Thinking Machines (2024).
\newblock Defeating nondeterminism in LLM inference.
\newblock {\em https://thinkingmachines.ai/blog/defeating-nondeterminism-in-llm-inference/}.

\bibitem{tang2023toolalpaca}
Tang, Q., et al. (2023).
\newblock ToolAlpaca: Generalized tool learning for language models with 3000 simulated cases.
\newblock {\em arXiv preprint arXiv:2306.05301}.

\bibitem{toucan2025}
beyoru (2025).
\newblock Toucan-1.5M: A Large-Scale Structured Tool-Calling Dataset.
\newblock {\em HuggingFace Datasets: https://huggingface.co/datasets/beyoru/Toucan-1.5M-structured-Qwen}.

\bibitem{wang2023selfconsistency}
Wang, X., et al. (2023).
\newblock Self-consistency improves chain of thought reasoning in language models.
\newblock {\em ICLR 2023}.

\bibitem{webson2022prompt}
Webson, A. and Pavlick, E. (2022).
\newblock Do prompt-based models really understand the meaning of their prompts?
\newblock {\em NAACL 2022}.

\bibitem{zhang1989simple}
Zhang, K. and Shasha, D. (1989).
\newblock Simple fast algorithms for the editing distance between trees and related problems.
\newblock {\em SIAM J. Computing}.

\bibitem{zhou2023large}
Zhou, Y., et al. (2023).
\newblock Large language models are human-level prompt engineers.
\newblock {\em ICLR 2023}.

\bibitem{ouyang2022training}
Ouyang, L., et al. (2022).
\newblock Training language models to follow instructions with human feedback.
\newblock {\em NeurIPS 2022}.

\bibitem{bai2022constitutional}
Bai, Y., et al. (2022).
\newblock Constitutional AI: Harmlessness from AI feedback.
\newblock {\em arXiv preprint arXiv:2212.08073}.

\bibitem{chen2023universal}
Chen, X., et al. (2023).
\newblock Universal self-consistency for large language model generation.
\newblock {\em arXiv preprint arXiv:2311.17311}.

\bibitem{zheng2024judging}
Zheng, L., et al. (2024).
\newblock Judging LLM-as-a-Judge with MT-Bench and Chatbot Arena.
\newblock {\em NeurIPS 2024}.

\bibitem{lundberg2017shap}
Lundberg, S. M. and Lee, S.-I. (2017).
\newblock A unified approach to interpreting model predictions.
\newblock {\em NeurIPS 2017}.

\bibitem{sharegpt2023}
Arun63 (2023).
\newblock ShareGPT Structured Output Datasets.
\newblock {\em HuggingFace Datasets: https://huggingface.co/datasets/Arun63/sharegpt-structured-output-json and https://huggingface.co/datasets/Arun63/sharegpt-quizz-generation-json-output}.

\end{thebibliography}

%==============================================================================
% Appendix
%==============================================================================
\appendix

\section{Weighted Analysis for Model Imbalance}
\label{app:weighted}

To address potential bias from Claude model overrepresentation (8 of 18 models, 44\%), we conduct weighted analysis where each model family receives equal aggregate weight through inverse frequency weighting.

\begin{table}[h]
\centering
\caption{Weighted vs Unweighted Analysis Comparison}
\label{tab:weighted}
\small
\begin{tabular}{@{}lccc@{}}
\toprule
\textbf{Metric} & \textbf{Unweighted} & \textbf{Weighted} & \textbf{$\Delta$} \\
\midrule
Schema Complexity $r$ & $-0.152$ & $-0.148$ & $<0.01$ \\
Query Length $r$ & $-0.069$ & $-0.071$ & $<0.01$ \\
Temperature $r$ & $-0.069$ & $-0.065$ & $<0.01$ \\
Model Rankings $\rho$ & --- & 0.94 & --- \\
Interaction Effect & 2.6$\times$ & 2.5$\times$ & Preserved \\
\bottomrule
\end{tabular}
\end{table}

Results demonstrate robustness: all correlations change by $<0.01$, model rankings remain highly stable (Spearman $\rho = 0.94$), and the temperature-complexity interaction effect (2.6$\times$ differential) is preserved under weighting.

\section{Sensitivity Analysis}
\label{app:sensitivity}

\textbf{Metric Sanity Check.} Table~\ref{tab:metric_sanity} validates both metrics on synthetic cases.

\begin{table}[h]
\centering
\caption{Metric Sanity Check on Synthetic Output Distributions}
\label{tab:metric_sanity}
\small
\begin{tabular}{@{}lcccc@{}}
\toprule
\textbf{Case} & \textbf{Similarities} & \textbf{$C_{mean}$} & \textbf{$S_\alpha$} & \textbf{Interpretation} \\
\midrule
Identical & all 1.0 & 1.00 & 1.00 & Best \\
High uniform & all 0.8 & 0.80 & 1.00 & Good, stable \\
Low uniform & all 0.2 & 0.20 & 1.00 & Poor but stable$^\dagger$ \\
Two clusters & 0.9/0.1 mix & 0.50 & 0.20 & Bimodal \\
Random & uniform [0,1] & 0.50 & 0.41 & Unstable \\
\bottomrule
\end{tabular}
\vspace{0.3em}

\raggedright\footnotesize $^\dagger$$S_\alpha$ alone would miss this; $C_{mean}=0.20$ correctly flags poor consistency.
\end{table}

We analyze the effect of the $\alpha$ parameter in $S_\alpha = (1/(1 + 2\hat{\sigma}))^\alpha$. Testing $\alpha \in \{1, 2, 5, 10\}$, relative model rankings are stable (Spearman $\rho > 0.95$ across all pairs). We use $\alpha=1$ for interpretability. Factor correlations remain directionally consistent across $\alpha$ values.

\textbf{$C_{mean}$ vs $S_\alpha$ Correlation.} Across our 199K instances, $C_{mean}$ and $S_\alpha$ correlate at $r=0.72$, indicating they capture related but distinct aspects. $C_{mean}$ better discriminates at high consistency levels; $S_\alpha$ better discriminates at low consistency levels. Our factor analysis results are robust to metric choice: top-3 features (schema complexity, query length, temperature) remain unchanged when using $C_{mean}$ as target.

\section{Experimental Details}
\label{app:experimental_details}

\textbf{Models (18 total):} Claude family (3.5-Haiku, 3.5-Sonnet, 3.7-Sonnet, Haiku-4.5, Sonnet-4, Sonnet-4.5, Opus-4, Opus-4.5), GPT-4.1-Mini, Qwen (32B, 235B), Llama-3.3-70B, Gemini-2.5-Flash-Lite, Grok-4.1-Fast, GPT-OSS-120B, Nova-2-Lite, Minimax-M2, Mimo-V2-Flash. Our dataset includes 8 Claude variants (44\%); weighted analysis (Appendix~\ref{app:weighted}) confirms robustness.

\textbf{Factor Categories (Toucan):} Tool Schema (4): tool count, total parameters, max params/tool, required params. Query (16): length, word count, sentence count, avg word length, vocabulary diversity, punctuation, question marks, non-ASCII ratio, linguistic markers. Model (1): identity. Config (1): temperature.

\textbf{Factor Categories (ShareGPT):} JSON Schema (3): json\_keys, json\_depth, schema\_complexity. Query features and temperature are shared across both datasets.

\section{Causal Validation Experimental Details}
\label{app:causal_details}

\textbf{Dataset.} We use samples from the Toucan benchmark~\cite{toucan2025}, filtering for samples with valid tool schemas (non-empty descriptions). This yields 847 valid samples from which we draw intervention candidates.

\textbf{Features Tested and Intervention Counts.} We test four linguistic features using bidirectional interventions (ADD and REMOVE). Table~\ref{tab:intervention_counts} shows the detailed breakdown:

\begin{table}[h]
\centering
\caption{Intervention Experiment Sample Counts by Feature}
\label{tab:intervention_counts}
\small
\begin{tabular}{@{}lcccp{4cm}@{}}
\toprule
\textbf{Feature} & \textbf{ADD} & \textbf{REMOVE} & \textbf{Total} & \textbf{Rewriting Rule} \\
\midrule
has\_if & 675 & 417 & 1,092 & Add/remove conditional clauses \\
has\_must & 299 & 360 & 659 & ``must'' $\leftrightarrow$ ``can'' \\
has\_should & 280 & 261 & 541 & ``should'' $\leftrightarrow$ ``can'' \\
has\_please & 220 & 192 & 412 & Add/remove ``please'' \\
\midrule
\textbf{Total} & 1,474 & 1,230 & 2,704 & --- \\
\bottomrule
\end{tabular}
\end{table}

\textbf{Model Coverage.} The 2,704 experiments are distributed across 4 models:
\begin{itemize}
    \item Llama-3.3-70B: 776 samples (all 4 features)
    \item Qwen3-32B: 655 samples (all 4 features)
    \item Nova-2-Lite: 671 samples (all 4 features)
    \item Claude-Sonnet-4: 602 samples (all 4 features)
\end{itemize}

\textbf{Rewriting Approach.} We use deterministic rule-based rewriting (not LLM-based) to ensure: (1) reproducibility across experiments, (2) isolation of the specific feature's effect, and (3) minimal perturbation to other prompt characteristics. Specific rewriting rules:
\begin{itemize}
    \item \textbf{has\_if}: ADD appends ``If possible, provide detailed results''; REMOVE strips conditional clauses.
    \item \textbf{has\_must}: Substitutes ``must'' with ``can'' while preserving sentence structure.
    \item \textbf{has\_should}: Substitutes ``should'' with ``can'' while preserving sentence structure.
    \item \textbf{has\_please}: ADD prepends ``Please''; REMOVE strips ``please'' occurrences.
\end{itemize}

\textbf{Experimental Parameters.}
\begin{itemize}
    \item Models: Claude-Sonnet-4, Nova-2-Lite, Llama-3.3-70B, Qwen3-32B
    \item Temperatures: 0.0, 0.5, 1.0
    \item Runs per condition: 10 independent API calls
    \item Total experiments: 2,704 intervention experiments (1,474 ADD + 1,230 REMOVE)
\end{itemize}

\textbf{Consistency Metric.} For computational efficiency at scale (2,704 experiments $\times$ 10 runs), we use mean pairwise Jaccard similarity of tool names as a lightweight proxy. This choice is validated against full STED: Pearson $r = 0.89$ ($p < 0.001$), indicating strong agreement. We additionally validate key findings using $C_{mean}$ from STED on a random subset (n=500), confirming directional consistency of all ceiling effect patterns. The Jaccard-STED correlation justifies using Jaccard for large-scale screening while STED remains the primary metric in main text analyses.

\textbf{Cross-Fitting Protocol.} To prevent regression-to-mean artifacts when stratifying by baseline:
\begin{itemize}
    \item Runs 1--5: Used to estimate baseline consistency (for stratification into low/high groups)
    \item Runs 6--10: Used to measure post-intervention effects (independent from stratification)
\end{itemize}
This ensures that noise in baseline estimation cannot inflate apparent intervention effects.

\textbf{Statistical Analysis.} We use paired t-tests comparing consistency before vs. after intervention. For conditional analysis, we partition samples by baseline consistency threshold (0.8) and apply Welch's t-test for unequal variances. Effect sizes are reported with 95\% confidence intervals where applicable.

\section{ADD vs REMOVE Intervention Details}
\label{app:add_remove}

Table~\ref{tab:add_remove} provides detailed results for ADD and REMOVE interventions separately, enabling assessment of bidirectional consistency.

\begin{table}[h]
\centering
\caption{Bidirectional Intervention Effects: ADD vs REMOVE}
\label{tab:add_remove}
\small
\begin{tabular}{@{}llcccc@{}}
\toprule
\textbf{Feature} & \textbf{Type} & \textbf{n} & \textbf{$C_{orig}$} & \textbf{$\Delta C$} & \textbf{p} \\
\midrule
has\_if & ADD & 675 & 0.319 & +0.5\% & 0.636 \\
has\_if & REMOVE & 417 & 0.319 & +4.2\% & 0.003 \\
\midrule
has\_must & ADD & 299 & 0.763 & +3.4\% & 0.012 \\
has\_must & REMOVE & 360 & 0.793 & +0.9\% & 0.133 \\
\midrule
has\_should & ADD & 280 & 0.807 & +0.0\% & 0.981 \\
has\_should & REMOVE & 261 & 0.738 & +1.3\% & 0.152 \\
\midrule
has\_please & ADD & 220 & 0.127 & $-$3.7\% & 0.009 \\
has\_please & REMOVE & 192 & 0.115 & +2.1\% & 0.103 \\
\bottomrule
\end{tabular}
\end{table}

\textbf{Bidirectional Patterns.} Several features show opposite effects for ADD vs REMOVE:
\begin{itemize}
    \item \texttt{has\_please}: ADD $\Delta C = -3.7\%$ (p=0.009), REMOVE $\Delta C = +2.1\%$ (p=0.103)
    \item \texttt{has\_if}: ADD $\Delta C = +0.5\%$, REMOVE $\Delta C = +4.2\%$ (p=0.003)
    \item \texttt{has\_must}: ADD $\Delta C = +3.4\%$ (p=0.012), REMOVE $\Delta C = +0.9\%$
    \item \texttt{has\_should}: ADD $\Delta C = +0.0\%$, REMOVE $\Delta C = +1.3\%$
\end{itemize}

\textbf{Stratified ADD Effects.} Table~\ref{tab:stratified_add} shows ADD intervention effects stratified by baseline consistency, revealing larger improvements for difficult prompts.

\begin{table}[h]
\centering
\caption{Stratified ADD Intervention Effects by Baseline}
\label{tab:stratified_add}
\small
\begin{tabular}{@{}lcccc@{}}
\toprule
& \multicolumn{2}{c}{\textbf{Low Baseline ($<$0.8)}} & \multicolumn{2}{c}{\textbf{High Baseline ($\geq$0.8)}} \\
\cmidrule(lr){2-3} \cmidrule(lr){4-5}
\textbf{Feature} & \textbf{n} & \textbf{$\Delta C$} & \textbf{n} & \textbf{$\Delta C$} \\
\midrule
has\_must & 91 & $\mathbf{+15.5\%}$ & 208 & $-1.9\%$ \\
has\_should & 164 & $\mathbf{+9.7\%}$ & 377 & $-3.4\%$ \\
has\_if & 464 & $\mathbf{+5.7\%}$ & 211 & $-11.0\%$ \\
has\_please & 192 & $\mathbf{+0.5\%}$ & 28 & $-32.9\%$ \\
\bottomrule
\end{tabular}
\end{table}

\textbf{Key Observation.} For all four features, ADD interventions on low-baseline prompts yield improvements (+0.5\% to +15.5\%), while high-baseline prompts show negative effects ($-1.9\%$ to $-32.9\%$). This confirms the ceiling effect pattern: interventions help difficult prompts but can harm easy ones. The largest stratified effect occurs for the directive features (\texttt{has\_must}: +15.5\%, \texttt{has\_should}: +9.7\%) on low-baseline prompts.

\section{Causal Validation Generalization Arguments}
\label{app:causal}

Our causal validation uses 4 models (Claude-Sonnet-4, Nova-2-Lite, Llama-3.3-70B, Qwen3-32B), providing empirical evidence for cross-model generalization:

\textbf{Argument 1: Consistent Ceiling Effect Across Models.} All 4 features show the same ceiling effect pattern across all 4 tested models: improvements for low-baseline prompts, decreases for high-baseline prompts. This consistency across architectures (Claude, Nova, Llama, Qwen) suggests the ceiling effect is a fundamental property of LLM consistency, not model-specific.

\textbf{Argument 2: LOMO Consistency.} Our LOMO analysis shows mean $R^2 = -0.009$ for cross-model transfer with high variance ($\pm 0.18$), indicating variable transfer depending on model similarity. However, the \textit{intervention effect direction} (positive for difficult, negative for easy) transfers across all models, suggesting the targeted optimization strategy generalizes.

\textbf{Argument 3: Feature Surface Agnosticism.} Our rewriting targets surface-level features (modal verbs, politeness markers, conditionals) that are model-agnostic by construction. These features do not interact specifically with any model's architecture, training data, or decoding strategy.

\textbf{Argument 4: Confounding Structure.} The observational associations arise from confounding by task complexity: naturally short prompts describe simpler tasks, while complex tasks require detailed prompts. This confounding structure is independent of model choice---all models face the same input distribution. Thus, the failure of mechanical rewriting (which preserves task complexity while changing surface features) should generalize across models.

\textbf{Argument 5: Consistency Across Temperature Settings.} The ceiling effect is consistent across all tested temperatures (0.0, 0.5, 1.0). At T=0.0, interventions on low-baseline prompts yield $+8.2\%$ improvement on average; at T=1.0, the improvement is $+11.4\%$. This suggests that the ceiling effect is not an artifact of sampling randomness but reflects fundamental properties of how models process prompts with different baseline difficulty levels.

\textbf{Argument 6: Accuracy Confirms Ceiling Effect.} Beyond consistency, we observe the same ceiling effect pattern in accuracy. For difficult prompts, directive interventions improve accuracy by $+4.5\%$ to $+6.0\%$, while easy prompts show negligible changes ($-0.1\%$ to $-0.5\%$). This dual confirmation in both metrics strengthens confidence in the ceiling effect as a genuine phenomenon rather than a measurement artifact.

\section{LLM-as-Judge Validation Details}
\label{app:llm_judge}

\begin{table}[h]
\centering
\caption{LLM-as-Judge Self-Consistency Analysis}
\small
\begin{tabular}{@{}lcc@{}}
\toprule
\textbf{Metric} & \textbf{T=0.0} & \textbf{T=0.7} \\
\midrule
All runs agree & 100\% & 52\% \\
Avg consistency & 100\% & 83.6\% \\
Agreement w/ STED & 90\% & 64\% \\
\bottomrule
\end{tabular}
\end{table}

We conduct 5 repeated judgments per sample. At T=0.7, only 52\% produce identical judgments---the very inconsistency we aim to measure.

\section{Pairwise Transfer Matrix}
\label{app:transfer}

\begin{table}[h]
\centering
\caption{Pairwise Model Transfer Matrix ($R^2$)}
\small
\begin{tabular}{@{}lcccc@{}}
\toprule
\textbf{Train$\backslash$Test} & \textbf{Nova} & \textbf{Sonnet-4} & \textbf{Opus-4.5} & \textbf{GPT-4.1} \\
\midrule
Nova-2-Lite & (0.34) & 0.10 & 0.12 & $-0.17$ \\
Claude-Sonnet-4 & 0.08 & (0.38) & 0.15 & $-0.20$ \\
Claude-Opus-4.5 & 0.05 & 0.09 & (0.34) & $-0.12$ \\
GPT-4.1-Mini & $-0.54$ & $-0.47$ & $-0.25$ & (0.38) \\
\bottomrule
\end{tabular}
\end{table}

GPT-4.1-Mini shows strongly negative transfer to all other models, explaining its outlier status.

\section{Per-Model SHAP Analysis (Full)}
\label{app:permodel_shap}

\begin{table}[h]
\centering
\caption{Per-Model SHAP Analysis: All 18 Models (Full Version)}
\label{tab:permodel_shap_full}
\small
\begin{tabular}{@{}lccc@{}}
\toprule
\textbf{Model} & \textbf{Top SHAP Feature} & \textbf{SHAP Value} & \textbf{Avg $\rho$} \\
\midrule
GPT-4.1-Mini & \textbf{schema\_depth} & 0.133 & \textbf{0.07} \\
Llama-3.3-70B & \textbf{schema\_breadth} & 0.112 & 0.53 \\
Gemini-2.5-Flash-Lite & query\_length & 0.133 & 0.66 \\
Minimax-M2 & \textbf{num\_tools} & 0.092 & 0.56 \\
GPT-OSS-120B & \textbf{avg\_params\_per\_tool} & 0.059 & 0.62 \\
Grok-4.1-Fast & schema\_complexity & 0.069 & 0.55 \\
Qwen3-235B-A22B & schema\_complexity & 0.070 & 0.63 \\
Qwen3-32B & query\_length & 0.046 & 0.69 \\
Claude-Opus-4 & schema\_complexity & 0.040 & 0.66 \\
Claude-Opus-4.5 & schema\_breadth & 0.035 & 0.58 \\
Claude-Sonnet-4 & avg\_tool\_name\_length & 0.038 & 0.65 \\
Claude-Sonnet-4.5 & query\_word\_count & 0.040 & 0.65 \\
Claude-3.7-Sonnet & query\_length & 0.038 & 0.68 \\
Claude-3.5-Sonnet & query\_length & 0.024 & 0.68 \\
Claude-Haiku-4.5 & schema\_complexity & 0.057 & 0.63 \\
Claude-3.5-Haiku & query\_length & 0.038 & 0.60 \\
Mimo-V2-Flash & tool\_name\_ambiguity & 0.063 & 0.66 \\
Nova-2-Lite & avg\_tool\_name\_length & 0.038 & 0.65 \\
\bottomrule
\end{tabular}
\end{table}

\textbf{Key Findings:} GPT-4.1-Mini ($\rho=0.07$) is a major outlier with completely different feature sensitivity. GPT-OSS-120B shows unique sensitivity to \texttt{avg\_params\_per\_tool}. Most Claude and Qwen models prioritize \texttt{query\_length} or \texttt{schema\_complexity}. SHAP vs MDI correlation across all models: mean $\rho=0.89$, range 0.86--0.96.

\section{Cross-Dataset Validation: ShareGPT}
\label{app:sharegpt}

We validate our findings on ShareGPT, a dataset of 80 JSON generation samples with explicit output schemas (15,840 evaluation instances across the same 18 models).

\subsection{Summary Comparison}

Table~\ref{tab:cross_dataset} compares key metrics between Toucan and ShareGPT.

\begin{table}[h]
\centering
\caption{Cross-Dataset Validation: Toucan vs ShareGPT}
\label{tab:cross_dataset}
\small
\begin{tabular}{@{}lcc@{}}
\toprule
\textbf{Metric} & \textbf{Toucan} & \textbf{ShareGPT} \\
\midrule
Samples & 1,006 & 80 \\
Evaluation instances & 199,188 & 14,960 \\
\midrule
Per-model $R^2$ (mean) & 0.67 & 0.60 \\
Pooled $R^2$ & 0.10 & 0.07 \\
$R^2$ Gap & 0.57 & 0.53 \\
\midrule
LOMO $R^2$ & $-0.009 \pm 0.18$ & $\approx 0$ \\
\bottomrule
\end{tabular}
\end{table}

\subsection{Predictive Model Comparison (ShareGPT)}

\begin{table}[h]
\centering
\caption{ShareGPT: Predictive Model Comparison (5-fold CV)}
\label{tab:sharegpt_model_comparison}
\small
\begin{tabular}{@{}lcc@{}}
\toprule
\textbf{Model} & \textbf{$R^2$} & \textbf{Std} \\
\midrule
Gradient Boosting & 0.102 & 0.097 \\
Random Forest & 0.074 & 0.100 \\
Ridge Regression & 0.050 & 0.072 \\
\bottomrule
\end{tabular}
\end{table}

\subsection{Per-Model $R^2$ (ShareGPT)}

\begin{table}[h]
\centering
\caption{ShareGPT: Per-Model $R^2$ from Controllable Factors (RF, 5-fold CV)}
\label{tab:sharegpt_permodel_r2}
\small
\begin{tabular}{@{}lcc@{}}
\toprule
\textbf{Model} & \textbf{$R^2$} & \textbf{Top Feature} \\
\midrule
Claude-3.5-Sonnet & 0.859 & json\_keys \\
Llama-3.3-70B & 0.791 & schema\_cplx \\
Mimo-V2-Flash & 0.753 & schema\_cplx \\
Qwen3-235B-A22B & 0.740 & schema\_cplx \\
Claude-3.7-Sonnet & 0.714 & json\_keys \\
Qwen3-32B & 0.695 & schema\_keys \\
Nova-2-Lite & 0.672 & temperature \\
Claude-Opus-4.5 & 0.642 & temperature \\
GPT-4.1-Mini & 0.641 & temperature \\
Grok-4.1-Fast & 0.586 & json\_keys \\
Claude-Haiku-4.5 & 0.585 & temperature \\
Claude-Sonnet-4.5 & 0.577 & temperature \\
Minimax-M2 & 0.463 & temperature \\
Claude-Sonnet-4 & 0.422 & temperature \\
Claude-Opus-4 & 0.420 & temperature \\
Claude-3.5-Haiku & 0.396 & temperature \\
Gemini-2.5-Flash-Lite & 0.306 & temperature \\
\midrule
\textbf{Mean} & \textbf{0.604} & --- \\
\textbf{Pooled} & 0.074 & --- \\
\bottomrule
\end{tabular}
\end{table}

\subsection{Feature Importance (ShareGPT)}

\begin{table}[h]
\centering
\caption{ShareGPT: Feature Importance (Pooled, MDI)}
\label{tab:sharegpt_feature_importance}
\small
\begin{tabular}{@{}lcc@{}}
\toprule
\textbf{Feature} & \textbf{Importance} & \textbf{Percentage} \\
\midrule
temperature & 0.348 & 34.8\% \\
json\_keys & 0.208 & 20.8\% \\
non\_ascii\_ratio & 0.132 & 13.2\% \\
avg\_word\_length & 0.101 & 10.1\% \\
json\_depth & 0.055 & 5.5\% \\
vocabulary\_diversity & 0.044 & 4.4\% \\
query\_complexity\_score & 0.036 & 3.6\% \\
\bottomrule
\end{tabular}
\end{table}

\textbf{Key difference from Toucan:} Temperature dominates ShareGPT (34.8\%) vs Toucan (11\%), since ShareGPT lacks tool schema features that compete for importance. Output complexity (\texttt{json\_keys}) replaces schema complexity as the second most important feature.

\textbf{Key Observations:} (1) The per-model vs pooled $R^2$ gap is nearly identical (0.57 vs 0.51), confirming that model heterogeneity causes effect cancellation across both datasets. (2) LOMO transfer is essentially zero in both cases, validating limited cross-model generalization. (3) Model family rankings are perfectly consistent---Claude leads in both datasets. (4) Temperature sensitivity differs: without tool schemas, temperature becomes the dominant factor (34.8\% vs 11\%).

\section{Temperature and Validity Analysis}
\label{app:temp_validity}

\begin{table}[h]
\centering
\caption{Temperature Effects on Validity vs Consistency}
\small
\begin{tabular}{@{}lccc@{}}
\toprule
\textbf{Temp} & \textbf{Validity} & \textbf{Consistency} & \textbf{$C_{mean}$} \\
\midrule
0.0 & 0.873 & 0.493 & 0.608 \\
0.3 & 0.873 & 0.458 & 0.615 \\
0.7 & 0.870 & 0.431 & 0.619 \\
1.0 & 0.871 & 0.409 & 0.622 \\
\bottomrule
\end{tabular}
\end{table}

\begin{table}[h]
\centering
\caption{Pareto-Optimal Models (High Validity + High Consistency)}
\small
\begin{tabular}{@{}lcc@{}}
\toprule
\textbf{Model} & \textbf{Validity} & \textbf{Consistency} \\
\midrule
Claude-Opus-4 & 0.971 & 0.638 \\
Claude-Sonnet-4 & 0.970 & 0.623 \\
Nova-2-Lite & 0.965 & 0.615 \\
Qwen3-235B-A22B & 0.981 & 0.529 \\
\bottomrule
\end{tabular}
\end{table}

\section{Confounding Analysis Details}
\label{app:confounding}

\begin{table}[h]
\centering
\caption{Baseline Consistency by Feature Presence}
\small
\begin{tabular}{@{}lcccc@{}}
\toprule
\textbf{Feature} & \textbf{Without} & \textbf{With} & \textbf{Diff.} & \textbf{p-value} \\
\midrule
has\_please & 0.876 & 0.966 & +0.091 & 0.0060 \\
has\_if & 0.919 & 0.924 & +0.005 & 0.855 \\
has\_must & 0.921 & 0.923 & +0.002 & 0.931 \\
has\_should & 0.920 & 0.925 & +0.005 & 0.823 \\
\bottomrule
\end{tabular}
\end{table}

\begin{table}[h]
\centering
\caption{Intervention Effect vs Baseline Consistency Correlation}
\small
\begin{tabular}{@{}lcc@{}}
\toprule
\textbf{Feature} & \textbf{Pearson $r$} & \textbf{p-value} \\
\midrule
has\_please & $-0.373$ & $<$0.0001 \\
has\_if & $-0.290$ & $<$0.0001 \\
has\_must & $-0.248$ & 0.0001 \\
has\_should & $-0.261$ & $<$0.0001 \\
\bottomrule
\end{tabular}
\end{table}

\section{Stability Score and Accuracy Tables}
\label{app:stability}

\begin{table}[h]
\centering
\caption{Stability Score ($S_\alpha$) Changes by Baseline Level}
\small
\begin{tabular}{@{}lcccccc@{}}
\toprule
& \multicolumn{2}{c}{\textbf{Low ($<$0.8)}} & \multicolumn{2}{c}{\textbf{High ($\geq$0.8)}} & & \\
\cmidrule(lr){2-3} \cmidrule(lr){4-5}
\textbf{Feature} & \textbf{$S_\alpha^{orig}$} & \textbf{$\Delta S_\alpha$} & \textbf{$S_\alpha^{orig}$} & \textbf{$\Delta S_\alpha$} & \textbf{t} & \textbf{p} \\
\midrule
has\_should & 0.021 & +0.368 & 0.973 & $-$0.067 & 8.29 & $<$0.0001 \\
has\_if & 0.038 & +0.304 & 0.956 & $-$0.060 & 6.81 & $<$0.0001 \\
has\_please & 0.026 & +0.259 & 0.983 & $-$0.080 & 6.61 & $<$0.0001 \\
\bottomrule
\end{tabular}
\end{table}

\begin{table}[h]
\centering
\caption{Stratified Accuracy by Baseline Consistency Level}
\small
\begin{tabular}{@{}lcccccc@{}}
\toprule
& \multicolumn{3}{c}{\textbf{Low Baseline}} & \multicolumn{3}{c}{\textbf{High Baseline}} \\
\cmidrule(lr){2-4} \cmidrule(lr){5-7}
\textbf{Feature} & \textbf{n} & \textbf{$Acc_{orig}$} & \textbf{$\Delta Acc$} & \textbf{n} & \textbf{$Acc_{orig}$} & \textbf{$\Delta Acc$} \\
\midrule
has\_should & 47 & 0.086 & +0.312 & 192 & 0.995 & $-$0.082 \\
has\_please & 41 & 0.222 & +0.174 & 233 & 0.997 & $-$0.096 \\
has\_if & 41 & 0.319 & +0.124 & 263 & 0.992 & $-$0.051 \\
\bottomrule
\end{tabular}
\end{table}

\section{VIF Analysis for Schema Variables}
\label{app:vif}

Table~\ref{tab:vif} shows the Variance Inflation Factor (VIF) for schema features. VIF measures multicollinearity: VIF $>$ 5 indicates moderate, VIF $>$ 10 indicates high multicollinearity.

\begin{table}[h]
\centering
\caption{VIF Analysis of Schema Features}
\label{tab:vif}
\small
\begin{tabular}{@{}lcc@{}}
\toprule
\textbf{Feature} & \textbf{VIF} & \textbf{Interpretation} \\
\midrule
schema\_complexity & 105.5 & High ($>$10) \\
schema\_depth & 58.9 & High ($>$10) \\
schema\_breadth & 41.8 & High ($>$10) \\
total\_params & 24.0 & High ($>$10) \\
tool\_prefix\_diversity & 7.1 & Moderate ($>$5) \\
max\_params\_per\_tool & 6.3 & Moderate ($>$5) \\
avg\_params\_per\_tool & 4.7 & OK \\
tool\_name\_ambiguity & 3.9 & OK \\
param\_type\_diversity & 3.2 & OK \\
avg\_tool\_name\_length & 2.0 & OK \\
num\_tools & 1.7 & OK \\
\bottomrule
\end{tabular}
\end{table}

\textbf{Highly correlated pairs} ($|r| > 0.7$):
\begin{itemize}
    \item schema\_complexity $\leftrightarrow$ schema\_depth: $r = 0.946$
    \item schema\_breadth $\leftrightarrow$ total\_params: $r = 0.955$
    \item schema\_complexity $\leftrightarrow$ schema\_breadth: $r = 0.813$
    \item max\_params\_per\_tool $\leftrightarrow$ avg\_params\_per\_tool: $r = 0.844$
\end{itemize}

High multicollinearity among schema features can cause worse CV performance when combined than separately, due to overfitting from redundant correlated features.

\section{Decision Point Hypothesis Test}
\label{app:decision_point}

We tested the hypothesis that consistency scales as $p^{D(s)}$ where $D(s)$ is the number of decision points (tool selections + parameter decisions).

\textbf{Method:} Estimate $D(s) = \text{num\_tools} + 2 \times \text{total\_params}$ and fit:
\begin{equation}
\log(\text{consistency}) = \alpha + \beta \cdot D(s) + \epsilon
\end{equation}

\textbf{Results:}
\begin{itemize}
    \item $R^2 < 0.001$ (negligible explained variance)
    \item $\beta = -0.00001$ (slope essentially zero)
    \item $p = 0.004$ (statistically significant but practically meaningless)
    \item Implied $p$ per decision $\approx 1.0000$
\end{itemize}

\textbf{Conclusion:} The multiplicative hypothesis is \textit{directionally correct} (negative correlation) but explains negligible variance. Model identity and schema structure dominate over simple decision counts. The hypothesis provides useful intuition but should not be used for quantitative prediction.

\section{Failure Mode Taxonomy}
\label{app:failure_modes}

Analysis of low-consistency samples ($S_\alpha < 0.5$) reveals three primary failure patterns:

\textbf{Tool Selection Instability (38\%).} The model selects different tools across runs for the same query, common with overlapping functionality or ambiguous names (e.g., \texttt{search\_products} vs \texttt{find\_items}).

\textbf{Parameter Omission Variance (35\%).} Inconsistent inclusion of optional parameters. Some runs include detailed optional fields while others omit them entirely.

\textbf{Value Format Variation (27\%).} Parameter values take different but semantically equivalent forms (e.g., ``2024-01-15'' vs ``January 15, 2024'', ``1000'' vs ``1,000'').

Structural mismatches are rare ($<$5\%)---models understand schema \textit{structure} but struggle with content-level \textit{decisions}.

\textbf{Implications for LLM Development:}
\begin{itemize}
    \item \textbf{Training-time}: Consistency as explicit training objective
    \item \textbf{Inference-time}: Ensemble methods or self-consistency checking
    \item \textbf{Tooling}: Automated schema analysis to flag high-inconsistency patterns
\end{itemize}

\end{document}
